\section{Hitting Times}

\begin{definition}[Hitting Time, Hitting Probability]
    Let \(A \subseteq I\). The \term{first hitting time} of \(A\), \(T_A\), is given by
    \[
        T_A = \inf\Bqty{n \geq 0 : X_n \in A},
    \]
    recalling \(\inf \emptyset = \plusinfty\).

    The \term{hitting probabilities} of \(A\) is given by \(\functiondomain{h^A}{I}{\ccintv{0}{1}}\) with
    \[
        h_i^A = \Prob_i \pqty{T_A < \infty}
    \]
    for \(i \in I\).

    The \term{mean hitting time} is \(\functiondomain{k^A}{I}{\ccintv{0}{\infty}}\), given by
    \[
        k_i^A = \Expt_i \pqty{T_A} = \sum_{n \geq 0} n \cdot \Prob_i \pqty{T_A = n} + \infty \cdot \Prob_i \pqty{T_A = \infty}.
    \]
\end{definition}

\begin{theorem}[Hitting Probabilities]
    \label{thm:hitting-probabilities}%
    Suppose \(X\) is a Markov chain on \(I\) with transition matrix \(P\). If \(A \subseteq I\), then the hitting probabilities \(h_i^A\) satisfy
    \[
        h_i^A = \begin{cases}
            1, & i \in A, \\
            \sum_{j \in I} P \pqty{i, j} h_j^A, & i \notin A.
        \end{cases}
    \]

    Moreover, \(\pqty{h_i}^A\) are the minimal non-negative solutions to the above system of equations.
\end{theorem}

\begin{proof}
    If \(i \in A\), then \(h_i^A = 1\) by definition. If \(i \notin A\), then
    \begin{align*}
        h_i^A &= \Prob_i \pqty{T_A < \infty} \\
        &= \Prob_i \pqty{\Bqty{T_A = 0} \cup \bigcup_{n \geq 1} \Bqty{T_A = n}} \\
        &= \sum_{n \geq 1} \Prob_i \pqty{T_A = n} \\
        &= \Prob\pqty{\Conditional{X_1 \in A}{X_0 = i}} + \sum_{n \geq 2} \Prob \pqty{\Conditional{X_1 \notin A, \ldots, X_{n-1} \notin A, X_n \in A}{X_0 = i}} \\
        &= \sum_{j \in I} P\pqty{i, j} \Prob\pqty{\Conditional{X_1 \in A}{X_0 = i, X_1 = j}} \\
        &\hspace{8em} + \sum_{n \geq 2} \sum_{j \in I} \Prob\pqty{\Conditional{X_1 \notin A, \ldots, X_{n-1} \notin A, X_n \in A}{X_0 = i, X_1 = j}} P\pqty{i, j}.
    \end{align*}

    Hence,
    \begin{align*}
        h_i^A &= \sum_{j \in I} P\pqty{i, j} \bqty{\Prob_j \pqty{X_0 \in A} + \sum_{n \geq 1} \Prob_j \pqty{X_0 \notin A, \ldots, X_{n-1} \notin A, X_n \in A}} \\
        &= \sum_{j \in I} P\pqty{i, j} \bqty{\Prob_j \pqty{T_A = 0} + \sum_{n \geq 1} \Prob_j\pqty{T_A = n}} \\
        &= \sum_{j \in I} P\pqty{i, j} \Prob_j \pqty{T_A < \infty} \\
        &= \sum_{j \in I} P\pqty{i, j} h_j^A.
    \end{align*}

    To show minimality, suppose \(h_i\) is another solution. If \(i \in A\), then \(h_i^A = h_i = 1\). If \(i \notin A\), we want to show \(h_i \geq h_i^A\). By assumption, we have
    \begin{align*}
        h_i &= \sum_{j \in I} P\pqty{i, j} h_j \\
        &= \sum_{j \in A} P\pqty{i, j} + \sum_{j \notin A} P\pqty{i, j} h_j \\
        &= \sum_{j \in A} P\pqty{i, j} + \sum_{j \notin A, k \in A} P\pqty{i, j} P\pqty{j, k} + \sum_{j \notin A, k \notin A} P\pqty{i, j} P\pqty{j, k} h_k.
    \end{align*}

    We iterate this to get
    \begin{align*}
        h_i &= \sum_{j_1 \in A} P\pqty{i, j_1} + \sum_{j_1 \notin A, j_2 \in A} P\pqty{i, j_1} P\pqty{j_1, j_2} + \cdots \\
        &\hspace{8em} + \sum_{j_1, \ldots, j_{n - 1} \notin A, j_{n} \in A} P\pqty{i, j_1} P\pqty{j_1, j_2} \cdots P\pqty{j_n, j_{n + 1}} \\
        &\hspace{8em} + \sum_{j_1, \ldots, j_n \notin A} P\pqty{i, j_1} P\pqty{j_1, j_2} \cdots P\pqty{j_n, k} h_k \\
        &\geq \sum_{j_1 \in A} P\pqty{i, j_1} + \sum_{j_1 \notin A, j_2 \in A} P\pqty{i, j_1} P\pqty{j_1, j_2} + \cdots \\
        &\hspace{8em} + \sum_{j_1, \ldots, j_{n - 1} \notin A, j_{n} \in A} P\pqty{i, j_1} P\pqty{j_1, j_2} \cdots P\pqty{j_n, j_{n + 1}} \\
        &= \Prob_i \pqty{X_1 \in A} + \Prob_i \pqty{X_1 \notin A, X_2 \in A} + \cdots + \Prob_i \pqty{X_1, \ldots, X_{n - 1} \notin A, X_{n} \in A} \\
        &= \Prob_i \pqty{T_A = 1} + \Prob_i \pqty{T_A = 2} + \cdots + \Prob_i \pqty{T_A = n} \\
        &= \Prob_i \pqty{T_A \leq n}.
    \end{align*}

    Therefore, taking \(n \to \infty\), \(h_i \geq \Prob_i \pqty{T_A < \infty} = h_i^A\).
\end{proof}

\begin{figure}[ht!]
    \centering
    \begin{tikzpicture}[->,>=Stealth,shorten >=1pt,auto,node distance=2cm,semithick,initial text=]
        \node[state] (q0)  {\(0\)};
        \node[state] (q1)  [right=of q0] {\(1\)};
        \node[state] (q2)  [right=of q1] {\(2\)};
        \node[state] (q3)  [right=of q2] {\(3\)};
        \path (q0) edge [loop left] node {\(1\)} (q0)
            (q1) edge node {\(1/2\)} (q0)
            (q1) edge [bend left=30] node {\(1/2\)} (q2)
            (q2) edge [bend left=30] node {\(1/2\)} (q1)
            (q2) edge node {\(1/2\)} (q3)
            (q3) edge [loop right] node {\(1\)} (q3);
    \end{tikzpicture}
    \caption{Hitting Probability Example}
    \label{fig:hitting-probability-example}%
\end{figure}

\begin{example}[Hitting Probability]
    Consider the Markov chain illustrated in \autoref{fig:hitting-probability-example}. Then the transition matrix is given by
    \[
        P = \pmqty{1 & 0 & 0 & 0 \\ 1/2 & 0 & 1/2 & 0 \\ 0 & 1/2 & 0 & 1/2 \\ 0 & 0 & 0 & 1}.
    \] 

    So there are two absorbing states, \(0\) and \(3\). We have \(h_3^A = 1\) and \(h_0^A = h_0^A\), with
    \begin{align*}
        h_1^A &= \frac{1}{2} h_0^A + \frac{1}{2} h_2^A \\
        h_2^A &= \frac{1}{2} h_1^A + \frac{1}{2} h_3^A.
    \end{align*}

    This is not enough information to solve for \(h_1^A\) and \(h_2^A\). However, we can use the minimal property and considering \(h_0^A = 0\) to get \(h_2^A = \flatfrac{2}{3}\) and \(h_1^A = \flatfrac{1}{3}\).

    Another way to find \(h_2^A\) is simply by evaluating an infinite geometric series
    \[
        h_2^A = \frac{1}{2} + \frac{1}{8} + \frac{1}{32} + \cdots = \frac{2}{3}
    \]
\end{example}

\begin{example}[Random Walk on Positive Integers]
    \label{ex:random-walk-hitting-probability}%
    Consider a random walk on \(\NN_0\) with reflection at the origin. Suppose \(X\) is a Markov chain on \(I = \NN_0\), with
    \[
        P\pqty{i, j} = \begin{cases}
            p, & j = i + 1, \\
            q = 1 - p, & j = i - 1, \\
            0, & \text{otherwise},
        \end{cases}
    \]
    for \(i \geq 1\), and \(P\pqty{0, 1} = 1\).

    We want to find \(h_i^{\Bqty{0}}\). Let \(A = \Bqty{0}\). From the theorem, we have \(h_0^A = 1\), and
    \[
        h_i^A = p h_{i + 1}^{A} + q h_{i - 1}^A, \quad i \geq 1.
    \]  

    \begin{itemize}
        \item If \(p \neq q\), we get the general solution as
        \[
            h_i^A = a + b \pqty{\frac{q}{p}}^i.
        \]

        \begin{itemize}
            \item If \(p > q\) (drifting towards infinity), then by minimality and non-negativity, we have \(a = 0\) and \(b = 1\). Then we have
            \[
                h_i^{A} = \pqty{\frac{q}{p}}^i.
            \]

            \item If \(p < q\) (drifting towards the origin), then minimality and non-negativity, we have \(a = 1\) and \(b = 0\), so \(h_i^A = 1\) for all \(i \geq 1\).
        \end{itemize}

        \item If \(p = q\), then the general solution of the form
        \[
            h_i^A = a + bi
        \]
        and again, minimality and non-negativity implies that \(a = 1\) and \(b = 0\), so \(h_i^A = 1\) for all \(i \geq 1\).
    \end{itemize}
\end{example}

\begin{example}[Birth and Death Chain on \(\NN_0\)]
    Suppose \(X\) is a Markov chain on \(\NN_0 = \Bqty{0, 1, \ldots}\), where the transition probabilities are given by
    \[
        P\pqty{i, i + 1} = p_i \qand P\pqty{i, i - 1} = 1 - p_i = q_i
    \]
    for the birth and death probabilities, and
    \[
        P\pqty{0, 0} = 1
    \]
    for the distinct state. What is the probability of distinction, i.e., \(h_i^A\) where \(A = \Bqty{0}\)?

    Clearly, if \(i = 0\), then \(h_i^A = 1\), and for \(i \geq 1\), the recursion equation
    \[
        h_i^A = p_i h_{i + 1}^A + q_i h_{i - 1}^A
    \]
    but note \(h_i^A = \pqty{p_i + q_i} h_i^A\), so
    \[
        p_i \pqty{h_{i + 1}^A - h_i^A} = q_i \pqty{h_i^A - h_{i - 1}^A}.
    \]

    Let \(u_i = h_i^A - h_{i - 1}^A\). Then we have
    \[
        u_{i + 1} = \frac{q_i}{p_i} u_i = \pqty{\prod_{j = 1}^{i} \frac{q_j}{p_j}} u_1
    \]
    and we let \(\gamma_i\) be the product for \(i \geq 1\), and \(\gamma_0 = 1\).

    So we have \(u_{i + 1} = \gamma_i u_1\). Therefore, by the definition of the \(u_i\)s, we have
    \begin{align*}
        h_i^A &= \sum_{j = 1}^{i} u_j + h_0^A \\
        &= \pqty{\sum_{j = 1}^{i} \gamma_{j - 1}} u_1 + 1.
    \end{align*}

    It remains to utilise minimality to obtain \(u_1\). Note that we want all probabilities to be non-negative, so noting \(u_1 = h_1^A - 1\).
    \[
        \pqty{h_1^A - 1} \sum_{j = 0}^{i - 1} \gamma_j + 1 \geq 0
    \]
    and hence
    \[
        h_1^A \geq 1 - \frac{1}{\sum_{j = 0}^{i - 1} \gamma_j}
    \]
    holds for all \(i\). Hence, since all the \(\gamma_j\)s are positive, we may conclude
    \[
        h_1^A \geq 1 - \frac{1}{\sum_{j = 0}^{\infty} \gamma_j}
    \]
    and denote this infinite sum as \(S\). So the minimality condition tells us that \(h_1^A = 1 - \flatfrac{1}{S}\), so \(u_1 = - \frac{1}{S}\).

    \begin{itemize}
        \item If \(S < \plusinfty\) is finite, then \(h_0^A = 1\) with \(h_1^A = 1 - \flatfrac{1}{S}\), with
        \[
            h_i^A = \frac{\sum_{j \geq i} \gamma_j}{\sum_{j \geq 0} \gamma_j}
        \]
        for all \(i \geq 1\).
        \item IF \(S = \plusinfty\) diverges, then \(h_0^A = 1\), \(h_1^A = 1\) and \(h_i^A = 1\) for all \(i\).
    \end{itemize}
\end{example}

Now we return to \autoref{fig:hitting-probability-example}.

\begin{example}[Mean Hitting Time]
    Refer to \autoref{fig:hitting-probability-example}, let \(A = \Bqty{0, 3}\), we have
    \[
        k_1^A = 1 + \frac{1}{2} k_2^A = 1 + \frac{1}{2} k_1^A
    \]
    so
    \[
        k_1^a = 2 = k_2^A.
    \]

    It is tempting to understand that the \(1\) arises from
    \[
        \frac{1}{2} + \frac{1}{2} \pqty{1 + k_2^A}
    \]
    while in fact, a better time to understand is from the expectation
    \begin{align*}
        \Expt\pqty{T} &= \sum_{n \geq 0} n \Prob\pqty{T = n} \\
        &= \sum_{n \geq 0} \Prob\pqty{T > n} \\
        &= 1 + \sum_{n \geq 1} \Prob\pqty{T > n},
    \end{align*}
    where \(1 = \Prob\pqty{T > 0}\).
\end{example}

\begin{theorem}
    Suppose \(X\) is a Markov chain on \(I\) with transition matrix \(P\), and let \(A \subseteq I\). Then the mean hitting times \(\pqty{k_i^A}_{i \in I}\) satisfy
    \[
        k_i^A = \begin{cases}
            0, & i \in A, \\
            1 + \sum_{j \in I} P\pqty{i, j} k_j^A, & i \notin A.
        \end{cases}
    \]

    Note this sum can be equivalently taken over \(j \notin A\). Moreover, \(\pqty{k_i^A}\) are the minimal non-negative solution to the system of equations.
\end{theorem}

\begin{proof}
    Clearly, \(k_i^A = 0\) if \(i \in A\) by definition. Let \(i \notin A\), then
    \begin{align*}
        k_i^A &= \Expt_i \pqty{T_A} \\
        &= \sum_{n \geq 0} \Prob_i \pqty{T_A > n} \\
        &= 1 + \sum_{n \geq 1} \Prob_i\pqty{T_A > n} \\
        &= 1 + \sum_{n \geq 1} \Prob\pqty{\Conditional{X_0 \notin A, X_1 \notin A, \ldots, X_n \notin A}{X_0 = i}} \\
        &= 1 + \sum_{n \geq 1} \sum_{j \in I} P\pqty{i, j} \Prob\pqty{\Conditional{X_1 \notin A, \ldots, X_n \notin A}{X_0 = i, X_1 = j}} \\
        &= 1 + \sum_{j \in I} P\pqty{i, j} \sum_{n \geq 1} \Prob\pqty{\Conditional{X_1 \notin A, \ldots, X_n \notin A}{X_1 = j}} \\
        &= 1 + \sum_{j \in I} P\pqty{i, j} \sum_{n \geq 0} \Prob\pqty{\Conditional{X_0 \notin A, \ldots, X_n \in A}{X_0 = j}} \\
        &= 1 + \sum_{j \in I} P\pqty{i, j} \sum_{n \geq 0} \Prob_j \pqty{T_A \geq n} \\
        &= 1 + \sum_{j \in I} P\pqty{i, j} k_j^A.
    \end{align*}

    To show minimality, suppose \(\pqty{k_i}\) is another set of solutions to the system of equations. Then, if \(i \in A\), \(k_i = 1 = k_i^A\), so let \(i \notin A\), then
    \begin{align*}
        k_i &= 1 + \sum_{j_1 \notin A} P \pqty{i, j_1} k_{j_1} \\
        &= 1 + \sum_{j_1 \notin A} P \pqty{i, j_1} + \sum_{j_1 \notin A, j_2 \notin A} P\pqty{i, j_1} P\pqty{j_1, j_2} + \cdots \\
        &\hspace{8em} + \sum_{j_1 \notin A, \ldots, j_n \notin A} P\pqty{i, j_1} P\pqty{j_1, j_2} \cdots P\pqty{j_{n - 1}, j_n} \\
        &\hspace{8em} + \sum_{j_1 \notin A, \ldots, j_{n + 1} \notin A} P\pqty{i, j_1} P\pqty{j_1, j_2} \cdots P\pqty{j_n, j_{n + 1}} k_{j_{n + 1}}.
    \end{align*}

    So
    \begin{align*}
        k_i &\geq 1 + \Prob_i \pqty{x_1 \notin A} + \Prob_i \pqty{x_1 \notin A, x_2 \notin A} + \cdots \\
        &\hspace{8em} + \Prob_i \pqty{X_1 \notin A, \ldots, X_n \notin A} \\
        &= 1 + \Prob_i \pqty{T_A > 1} + \Prob_i \pqty{T_A > 2} + \cdots + \Prob_i \pqty{T_A > n} \\
        &= \Prob_i \pqty{T_A > 0} + \Prob_i \pqty{T_A > 1} + \Prob_i \pqty{T_A > 2} + \cdots + \Prob_i \pqty{T_A > n},
    \end{align*}
    so for all \(n \in \NN\), we have
    \[
        k_i \geq \sum_{k = 0}^{n} \Prob_i\pqty{T_A > k}
    \]
    so take \(n \to \infty\) gives \(k_i \geq k_i^A\) finishing the proof.
\end{proof}